\documentclass[a4paper,12pt]{article}
\usepackage[utf8]{inputenc}
\usepackage[T1]{fontenc}
\usepackage[ngerman]{babel}
\usepackage{amsmath}
\usepackage{amssymb}
\usepackage{enumitem}
\usepackage{geometry}
\geometry{a4paper, margin=2.5cm}
\usepackage{parskip}

\title{Einsendeaufgaben -- Aufgabe 4.2\\[0.5em]
\large Modul 61112: Lineare Algebra}
\author{}
\date{}

\begin{document}
\maketitle

\section*{Aufgabe 4.2}

Wir ziehen die erste Zeile mit $-1$ als Vielfaches von den anderen Zeilen ab:
\[
\begin{pmatrix}
1+x_1 & x_2   & x_3   & \cdots & x_{n-1} & x_n \\
x_1   & 1+x_2 & x_3   & \cdots & x_{n-1} & x_n \\
x_1   & x_2   & 1+x_3 & \cdots & x_{n-1} & x_n \\
\vdots & \vdots & \vdots & \ddots & \vdots & \vdots \\
x_1   & x_2   & x_3   & \cdots & 1+x_{n-1} & x_n \\
x_1   & x_2   & x_3   & \cdots & x_{n-1} & 1+x_n
\end{pmatrix}
\]
\[
\rightsquigarrow
\]
\[
\begin{pmatrix}
1+x_1 & x_2 & x_3 & \cdots & x_{n-1} & x_n \\
-1    & 1   & 0   & \cdots & 0       & 0 \\
-1    & 0   & 1   & \cdots & 0       & 0 \\
\vdots & \vdots & \vdots & \ddots & \vdots & \vdots \\
-1    & 0   & 0   & \cdots & 1       & 0 \\
-1    & 0   & 0   & \cdots & 0       & 1
\end{pmatrix}
\]

Wir nennen die transformierte Matrix $A'_n$. Nach Satz 4.2.27(3) gilt $\det(A_n) = \det(A'_n)$. Wir werden also $\det(A'_n) = 1 + \sum_{k=1}^{n} x_k$ zeigen. Daraus folgt dann $\det(A_n) = 1 + \sum_{k=1}^{n} x_k$.

Wir zeigen $\det(A'_n) = 1 + \sum_{k=1}^{n} x_k$ mithilfe von Vollständiger Induktion:

\textbf{Induktionsanfang:} $n=1$

Es folgt direkt:
\[
\det(A'_1) = 1 + x_1 = 1 + \sum_{k=1}^{1} x_k
\]

\textbf{Induktionsschritt:} $n \to n+1$

Wir untersuchen $A_{n+1}$:
\[
\begin{pmatrix}
1+x_1 & x_2 & x_3 & \cdots & x_{n-1} & x_n & x_{n+1} \\
-1    & 1   & 0   & \cdots & 0       & 0   & 0 \\
-1    & 0   & 1   & \cdots & 0       & 0   & 0 \\
\vdots & \vdots & \vdots & \ddots & \vdots & \vdots & \vdots \\
-1    & 0   & 0   & \cdots & 1       & 0   & 0 \\
-1    & 0   & 0   & \cdots & 0       & 1   & 0 \\
-1    & 0   & 0   & \cdots & 0       & 0   & 1
\end{pmatrix}
\]

Wir entwickeln $A_{n+1}$ an der $n{+}1$-ten Spalte nach dem laplace'schen Entwicklungssatz, sodass folgt
\[
\det(A'_{n+1}) = \det(A'_{n+1,n+1}) + (-1)^{n+2} x_{n+1} \det(A'_{1,n+1}).
\]

Wir untersuchen zunächst die Unterdeterminante $A'_{1,n+1}$. Wir vertauschen die letzte Zeile mit vorletzten, dann mit der vorvorletzten usw., bis die letzte Zeile von $A'_{1,n+1}$ oben ist:
\[
A'_{1,n+1} =
\begin{pmatrix}
-1 & 1 & 0 & \cdots & 0 \\
-1 & 0 & 1 & \cdots & 0 \\
\vdots & \vdots & \vdots & \ddots & \vdots \\
-1 & 0 & 0 & \cdots & 1 \\
-1 & 0 & 0 & \cdots & 0
\end{pmatrix}
\]
\[
\rightsquigarrow
\begin{pmatrix}
-1 & 1 & 0 & \cdots & 0 \\
-1 & 0 & 1 & \cdots & 0 \\
\vdots & \vdots & \vdots & \ddots & \vdots \\
-1 & 0 & 0 & \cdots & 0 \\
-1 & 0 & 0 & \cdots & 1
\end{pmatrix}
\]
\[
\rightsquigarrow
\begin{pmatrix}
-1 & 1 & 0 & \cdots & 0 & 0 \\
-1 & 0 & 0 & \cdots & 0 & 0 \\
\vdots & \vdots & \vdots & \ddots & \vdots & \vdots \\
-1 & 0 & 0 & \cdots & 1 & 0 \\
-1 & 0 & 0 & \cdots & 0 & 1
\end{pmatrix}
\]
\[
\rightsquigarrow
\begin{pmatrix}
-1 & 0 & 0 & \cdots & 0 & 0 \\
-1 & 1 & 0 & \cdots & 0 & 0 \\
\vdots & \vdots & \vdots & \ddots & \vdots & \vdots \\
-1 & 0 & 0 & \cdots & 1 & 0 \\
-1 & 0 & 0 & \cdots & 0 & 1
\end{pmatrix}
\]

Wir nennen das neue Minor $A''_{1,n+1}$. Da $A''$ eine Dreiecksmatrix ist, folgt
\[
\det(A''_{1,n+1}) = \prod_{k=1}^{n} a_{kk} = -1.
\]

Da $A''_{1,n+1}$ durch $n-1$ Vertauschungen aus $A'_{1,n+1}$ konstruiert wurde, folgt
\begin{align*}
\det(A'_{1,n+1}) &= (-1)^{n-1} \det(A''_{1,n+1}) \\
                 &= (-1)^{n-1} (-1) \\
                 &= (-1)^{n}.
\end{align*}

Daraus können wir $\det(A'_{n+1})$ ermitteln:
\begin{align*}
\det(A'_{n+1}) &= \det(A'_{n+1,n+1}) + (-1)^{n+2} x_{n+1} \det(A'_{1,n+1}) \\
               &= \det(A'_{n+1,n+1}) + (-1)^{n+2} x_{n+1} (-1)^{n} \\
               &= \det(A'_{n+1,n+1}) + x_{n+1} \\
               &= 1 + \sum_{k=1}^{n} x_k + x_{n+1} && (\text{Induktionsvoraussetzung}) \\
               &= 1 + \sum_{k=1}^{n+1} x_k
\end{align*}

\end{document}
